
Aus dem zuvor beschriebenen Digitalisierungsdruck unter Ressourcenmangel
ergibt sich für KMU die Herausforderung, digitale Technologien so einzusetzen,
dass vorhandene finanzielle, personelle und organisatorische Ressourcen
möglichst wirksam genutzt werden. ERP-Systeme können in diesem Zusammenhang
als Enabler wirken, da sie zentrale Geschäftsprozesse in einem gemeinsamen
System zusammenführen und dadurch eine einheitliche Daten- und Prozessbasis
schaffen \parencite{Haddara2011,Nach2008}.

Gerade bei knappen personellen Ressourcen kann die Integration verschiedener
Unternehmensbereiche einen wichtigen Beitrag leisten. Informationen müssen
nicht mehr mehrfach in unterschiedlichen Systemen oder Dateien gepflegt werden,
sondern können zentral bereitgestellt und unternehmensweit genutzt werden \parencite{Tereshchenko2016,Leyh2017}.
Dadurch lassen sich Doppelarbeiten, Medienbrüche und Abstimmungsaufwände
reduzieren \parencite{Tereshchenko2016,Leyh2017}. 
Für KMU bedeutet dies, dass vorhandene Mitarbeitende weniger Zeit
für manuelle Koordination und Informationssuche aufwenden müssen.

ERP-Systeme können außerdem dabei helfen, mit steigenden Anforderungen aus
globalisierten Märkten, Partnerschaften und Wertschöpfungsnetzwerken
umzugehen \parencite{Leyh2015}. Gerade weil in und zwischen KMU immer größere Informationsflüsse
entstehen, kann ein integriertes System dazu beitragen, Informationen
strukturierter zu erfassen, zu verarbeiten und verfügbar zu machen
\parencite{Haddara2011}. Damit kann ein ERP-System KMU unterstützen, trotz
begrenzter Ressourcen handlungsfähig zu bleiben und externe Anforderungen
besser zu koordinieren.

Auch standardisierte Abläufe und systemgestützte Workflows können helfen,
begrenzte Ressourcen effizienter einzusetzen \parencite{Leyh2015, Setiawati2014}. Wiederkehrende Tätigkeiten lassen
sich durch ERP-Systeme strukturieren, unterstützen oder teilweise
automatisieren \parencite{Leyh2015}. Dadurch können operative Aufgaben schneller bearbeitet werden,
ohne dass dafür zwingend zusätzliche personelle Kapazitäten aufgebaut werden
müssen \parencite{Leyh2015}. ERP-Systeme können somit dazu beitragen, die Leistungsfähigkeit eines
KMU trotz begrenzter Ressourcen zu erhöhen.

Darüber hinaus kann eine zentrale Datenbasis die Entscheidungsfähigkeit
verbessern. Studien im KMU-Kontext zeigen, dass ERP-Systeme einen positiven
Einfluss auf Sichtbarkeit, Qualität und Kontrolle von Informationen haben
können und dadurch Entscheidungsprozesse unterstützt werden
\parencite{Haddara2011}. 
Gerade für KMU, die häufig nicht über große Analyse-
oder Controlling-Abteilungen verfügen, kann dies eine wichtige Unterstützung
darstellen. ERP-Systeme schaffen damit eine technologische Grundlage, um
Digitalisierungsanforderungen trotz Ressourcenknappheit besser bewältigen zu
können.
